\section{Goal drift via sub-task injection: definition, questions and hypotheses}
\label{sec:definition}

\subsection{Operational definition}

Let a session start with an objective $\Gz$ given by the user at $t=0$. We assume that $\Gz$ can be
decomposed into a finite set of checkable requirements $\Req=\{r_1,\dots,r_m\}$, for instance hidden
acceptance tests, and that it comes with an \emph{admissible scope} $S(\Gz)$: the set of files or
artefacts that a competent and faithful implementation is expected to modify. The agent produces a
trajectory $\tau=(a_1,\dots,a_T)$ of actions (tool calls and messages) and leaves a final repository
state $\sigma_T$.

A \textbf{sub-task} $s_i$ is any objective not contained in $\Gz$ that enters the agent's context
at some time $t_i>0$, either as a message or when the agent first encounters it in the environment.
We characterise each sub-task along two dimensions.

\emph{Source} $\mathrm{src}(s_i)$, i.e.\ how the sub-task enters the context:
\begin{itemize}
  \item \textsf{user}: a side request from the user in the middle of the session;
  \item \textsf{env}: an obstacle in the environment, e.g.\ an unrelated failing test;
  \item \textsf{text}: instruction-like text found in the environment, e.g.\ a \texttt{TODO} or
        README note;
  \item \textsf{cascade}: an error that must be resolved for $\Gz$ to be reachable, e.g.\ a broken
        import or a buggy helper.
\end{itemize}
In our implementation (Section~\ref{sec:method}), \textsf{user} sub-tasks arrive as messages at
fixed pauses in the session, whereas \textsf{env}, \textsf{text} and \textsf{cascade} sub-tasks are
planted in the repository before the session starts and enter the context only when the agent runs
the tests or reads the affected files.

\emph{Legitimacy} $\kappa(s_i)$ with respect to the principal's intent:
\begin{itemize}
  \item \textsf{sanctioned}: the user asked for it;
  \item \textsf{necessary}: $\Gz$ cannot be completed without it;
  \item \textsf{distractor}: neither, though it may look reasonable.
\end{itemize}
Sanctioned and necessary sub-tasks legitimately widen the scope. Writing $S(s_i)$ for the files a
sub-task legitimately touches, the admissible scope after injection is
\begin{equation}
  S^\ast = S(\Gz)\;\cup\!\!\bigcup_{i:\,\kappa(s_i)\in\{\textsf{sanctioned},\textsf{necessary}\}}\!\! S(s_i).
\end{equation}

Following the distinction between drift through actions and drift through inaction (failing to
undo what a competing objective left behind) introduced by \citet{arike2025goaldrift}, and the alignment checks of \citet{jia2024taskshield} and
\citet{jia2025gpa}, we consider five families of indicators:
\begin{enumerate}
  \item \textbf{Outcome loss}: the fraction of $\Req$ satisfied by $\sigma_T$, written
        $\mathrm{pass}(\tau)$, together with its variance across repeated runs. Following
        \citet{laban2025lost}, we separate aptitude, the best-case performance, from unreliability,
        the spread between runs. Sanctioned \textsf{user} sub-tasks are scored separately by their
        own tests, so that an agent which ignores a legitimate request is not mistaken for a
        focused one.
  \item \textbf{Scope violation}: the out-of-scope ratio
        \[
          \mathrm{OOS}(\tau)=\frac{|\Delta\sigma\setminus S^\ast|}{|\Delta\sigma|},
        \]
        where $|\Delta\sigma|$ is the number of changed lines and $\Delta\sigma\setminus S^\ast$ are the
        changes outside the admissible scope. This includes any modification of protected artefacts
        such as existing tests, which we also count separately.
  \item \textbf{Effort share}: the fraction of actions that serve neither $\Gz$ nor a sanctioned or
        necessary sub-task, as judged per action by a validated LLM judge:
        \[
          \mathrm{EFF}(\tau)=\frac{1}{T}\sum_{t=1}^{T} \mathbb{1}\big[a_t \text{ serves no admissible objective}\big].
        \]
  \item \textbf{Representational drift}: whether the agent, probed on a forked copy of the session
        so that the main trajectory is not perturbed \citep{ding2026contextecho}, can still state
        $\Gz$ and its completion status (\emph{goal recall}). \emph{Goal replacement} is the extreme
        case in which the agent's own final summary presents a sub-task as the main accomplishment.
  \item \textbf{Report fidelity}: agreement between what the agent's final summary claims and what
        $\sigma_T$ and the tests show, including omitted out-of-scope changes and over-claimed
        requirements. This indicator links drift to the supervisor's view
        \citep{bansal2024communication,tang2026fail}.
\end{enumerate}

\paragraph{Drift as a counterfactual.} A low pass rate does not by itself indicate drift: the agent
may simply be incapable. We therefore define \textbf{goal drift via sub-task injection} as the
degradation of these indicators in a trajectory with injected sub-tasks, relative to a matched
trajectory on the same task without them. For a dose of $k$ injected sub-tasks from source
$\mathrm{src}$,
\begin{equation}
  D_k^{\mathrm{src}} \;=\; \mathbb{E}\big[\mathrm{pass}(\tau_0)\big]-\mathbb{E}\big[\mathrm{pass}(\tau_k^{\mathrm{src}})\big],
\end{equation}
and analogously for $\mathrm{OOS}$, $\mathrm{EFF}$ and recall. Every session has the same structure (an initial segment, a fixed number of pauses and a final
segment), so conditions differ only in what the pauses and the repository contain. Two controls are
essential:
\begin{itemize}
  \item a \emph{neutral-interruption} control, in which every pause in the session carries
        ``please continue'' instead of a sub-task; this is the $k=0$ trajectory $\tau_0$;
  \item a \emph{matched-length filler} control, in which pauses carry irrelevant status messages of
        similar length that request no work.
\end{itemize}
The difference between $\tau_k^{\mathrm{src}}$ and the filler condition isolates the effect of
\emph{new objectives} from the effect of a \emph{longer context}. Long-context studies vary the
amount, position or relevance of context \citep{liu2024lost,xue2026contextrot}, but not whether it
carries new objectives; that separation is required to attribute drift to sub-tasks.

\subsection{Research questions}

\begin{description}[leftmargin=0pt]
  \item[RQ1 (dose and source).] Does goal-level performance degrade as the number of injected
        sub-tasks grows, beyond the effect of context length alone? Do the four sources produce
        different drift signatures?
  \item[RQ2 (harness mitigations).] Which features of the agent harness reduce drift: re-injecting
        $\Gz$ through hooks, an external requirement checklist verified when the agent stops,
        delegating sub-tasks to sub-agents with isolated context, or an enforced task list?
  \item[RQ3 (compaction and inheritance).] Does compaction of the context, or resuming from a
        trajectory that has already drifted, transmit or amplify drift?
  \item[RQ4 (supervisor's view).] How faithfully do the agent's own reports reflect drift? How much
        drift can a supervisor intercept, and at what cost in interventions, under different
        checkpoint policies?
\end{description}

\subsection{Hypotheses}

\begin{description}[leftmargin=0pt]
  \item[H1 (dose--response).] $\mathrm{pass}$ decreases and $\mathrm{OOS}$ and $\mathrm{EFF}$ increase
        monotonically with $k$. The effect is larger than in the matched-length filler control,
        consistent with the dilution and pattern-matching mechanisms
        \citep{arike2025goaldrift,xue2026contextrot}.
  \item[H2 (source signatures).] Sources differ in \emph{how} they produce drift:
        \begin{itemize}
          \item sanctioned \textsf{user} requests mainly cause \emph{incompletion}, i.e.\ the agent
                does not return to $\Gz$ \citep{laban2025lost,zou2026interruptbench};
          \item \textsf{text} distractors mainly cause \emph{scope violation}, because the agent
                follows instruction-shaped artefacts \citep{mavali2026tab,ma2024caution};
          \item \textsf{env} obstacles mainly cause \emph{effort displacement} and changes to
                protected tests \citep{zhong2025impossiblebench,gloaguen2026fixedbench};
          \item necessary \textsf{cascade} sub-tasks cause the least drift per unit of effort, because
                resolving them serves $\Gz$.
        \end{itemize}
  \item[H3 (mitigations).] Mechanisms that act outside the model's context reduce drift more than
        instructions written into it:
        \begin{itemize}
          \item deterministic hooks that re-inject $\Gz$ or block completion until $\Req$ is checked;
          \item context isolation through sub-agents \citep{fei2026codedelegator}.
        \end{itemize}
        Merely asking the agent to ``stay focused'' or to keep a task list does less, because a task
        list can absorb distractors as easily as it anchors the goal.
  \item[H4 (compaction and inheritance).] Forced compaction lowers goal recall and increases constraint
        violations after the compaction point. Resuming from a drifted prefix produced by a weaker
        model transfers drift to a stronger one \citep{menon2026inherited}.
  \item[H5 (observability).] Agents' final reports under-report out-of-scope work and over-report
        completion, so report fidelity is below 1 and decreases with $k$ \citep{tang2026fail}.
        Drift-triggered checkpoints intercept more drift per human intervention than fixed
        per-step reviews, but the monitors that trigger them degrade on longer trajectories
        \citep{martin2026classifier}.
\end{description}

H1--H4 concern the agent as an operator; H5 concerns the human as its supervisor. The experimental
programme in the following sections is designed so that the same trajectories serve both levels of
analysis.
